\documentclass[11pt]{amsart}
\usepackage{amssymb,amsmath,amsthm,mathtools}
\usepackage[margin=1.1in]{geometry}
\usepackage{enumitem}
\usepackage[hidelinks]{hyperref}

\newtheorem{theorem}{Theorem}
\newtheorem{proposition}[theorem]{Proposition}
\newtheorem{lemma}[theorem]{Lemma}
\newtheorem{corollary}[theorem]{Corollary}
\theoremstyle{definition}
\newtheorem{definition}[theorem]{Definition}
\newtheorem{setup}[theorem]{Setup}
\theoremstyle{remark}
\newtheorem{remark}[theorem]{Remark}

\newcommand{\QQ}{\mathbb{Q}}
\newcommand{\CC}{\mathbb{C}}
\newcommand{\Hg}{\operatorname{Hg}}
\newcommand{\SU}{\operatorname{SU}}
\newcommand{\Sh}{\mathrm{Sh}}
\newcommand{\Hilb}{\operatorname{Hilb}}
\newcommand{\cl}{\operatorname{cl}}

\title[Closing the algebraic locus of Weil classes]{Closing the algebraic locus of Weil classes:\\
an attempt and the exact point where it stops}
\author{Deep Bhattacharjee}
\date{October 2026}

\begin{document}

\begin{abstract}
On a general abelian variety of Weil type the Hodge conjecture comes down
to one question: is the Weil class algebraic? The set of members of a Weil
family on which it is algebraic is known to be dense, and to be either the
whole family or a meagre set. This note tries to show that it is the whole
family, by the most direct route: take a member where the class is
algebraic and move the cycle to every other member. We prove what that
route gives. A stratum of the algebraic locus that passes through a member
with the largest possible Hodge group has the largest possible monodromy.
A single semiregular cycle at any one member would close the locus. We then
show where the route stops. The only members at which the class is known to
be algebraic have small Hodge group, and no semiregular cycle with the
right class is known. The attempt does \emph{not} prove the Hodge
conjecture for these varieties, and we say precisely what is missing.
\end{abstract}

\maketitle

\section{What we are trying to prove}

Fix an imaginary quadratic field $K$ and an integer $n\ge2$. An abelian
variety $A$ of dimension $2n$ is \emph{of Weil type} for $K$ if $K$ acts on
it by endomorphisms (up to isogeny), the Rosati involution of a
polarisation $\eta$ acts on $K$ as complex conjugation, and the two
eigenspaces of $K$ on the tangent space of $A$ both have dimension $n$.
Then $V=H^{1}(A,\QQ)$ is a $K$-vector space of dimension $2n$, the
polarisation gives a $K$-hermitian form $H$ on it of signature $(n,n)$, and
the $K$-determinant
\[
  \textstyle\bigwedge^{2n}_{K}V\subset\bigwedge^{2n}_{\QQ}V=H^{2n}(A,\QQ)
\]
is a plane of Hodge classes, the \emph{Weil classes} \cite{Weil77}. Weil
showed that on a general such $A$ they are not products of divisor
classes. In fact more is true: if the Hodge group of $A$ is the whole
special unitary group $\SU(V,H)$, which happens for a very general member
of the family, then the Hodge ring of $A$ is spanned by the powers of
$\eta$ together with this plane, and the Hodge conjecture for $A$ in every
degree is equivalent to the algebraicity of a single Weil class $\omega$
\cite{MZ99,BhH8}.

These varieties sit in one family. Let $D$ be the bounded symmetric domain
of $\SU(n,n)$, of dimension $n^{2}$, and let $\Gamma$ be a torsion-free
arithmetic subgroup of $\SU(V,H)(\QQ)$. Then $\Sh=\Gamma\backslash D$ is a
smooth quasi-projective variety (Baily and Borel \cite{BB66}), and over it
there is a projective abelian scheme $\pi\colon\mathcal A\to\Sh$ whose
fibres are the abelian varieties of Weil type with the given discriminant.
Since $\Gamma$ lies in the special unitary group, it acts trivially on the
$K$-determinant. So the Weil classes form a \emph{constant} local system
of rank two inside $R^{2n}\pi_{*}\QQ$, and $\omega$ is a global section of
it.

\begin{definition}
The \emph{algebraic locus} is the set $\Sigma\subseteq\Sh$ of points $t$
at which $\omega(t)$ is the class of an algebraic cycle with rational
coefficients on $A_{t}$.
\end{definition}

The goal is $\Sigma=\Sh$. That statement, for every $K$, $n$ and
discriminant, is the Hodge conjecture for every abelian variety of Weil
type. For $n=2$ it is a theorem of Markman \cite{Mar25a,Mar25b}. For
$n=3$ and a split hermitian form it is also Markman's theorem. For $n\ge3$
in general it is open.

\section{What is already known about the locus}

The first three facts are proved in \cite{BhH8}. We include proofs of the
first two, because they are short and the rest of the note uses the
argument.

\begin{lemma}\label{lem:countable}
$\Sigma$ is a countable union of closed algebraic subsets of $\Sh$.
\end{lemma}

\begin{proof}
Fix a relatively ample line bundle on $\mathcal A/\Sh$. For each Hilbert
polynomial $P$ the relative Hilbert scheme $\Hilb_{P}(\mathcal A/\Sh)$ is
projective over $\Sh$ and has finitely many irreducible components. Take
finitely many components $H_{1},\dots,H_{r}$, possibly from different
Hilbert schemes, and rational numbers $q_{1},\dots,q_{r}$. Over the fibre
product $H=H_{1}\times_{\Sh}\dots\times_{\Sh}H_{r}$ we have the cohomology
class $\sum_{i}q_{i}[Z_{i}]$ in the fibre $H^{2n}(A_{t},\QQ)$. Classes of
the fibres of a flat family of subschemes are locally constant. So on each
connected component of $H$ the class $\sum q_{i}[Z_{i}]$ either equals
$\omega(t)$ everywhere or nowhere. The union of the components where it
equals $\omega(t)$ maps properly to $\Sh$, so its image is closed and
algebraic. There are countably many choices of $P$'s, components and
$q_{i}$'s, and every point of $\Sigma$ lies in one of these images.
\end{proof}

\begin{lemma}[Dichotomy]\label{lem:dichotomy}
Either $\Sigma=\Sh$, or $\Sigma$ is meagre: it lies in a countable union
of closed subsets with empty interior, and its complement is dense.
\end{lemma}

\begin{proof}
$\Sh$ is irreducible. If one of the closed sets of
Lemma~\ref{lem:countable} is all of $\Sh$, then $\Sigma=\Sh$. Otherwise
each of them is a proper closed algebraic subset, which has empty
interior, and Baire's theorem on the complex manifold $\Sh$ gives the rest.
\end{proof}

\begin{proposition}[\cite{BhH8}]\label{prop:dense}
$\Sigma$ contains an explicit point $s_{0}$ at which $A_{s_{0}}$ is
isogenous to a product of abelian varieties with complex multiplication,
and at which $\omega$ is a polynomial in divisor classes. Its orbit under
the Hecke correspondences lies in $\Sigma$ and is dense in $\Sh$.
\end{proposition}

So $\Sigma$ is dense, and the problem is entirely about which side of the
dichotomy we are on. One stratum of Lemma~\ref{lem:countable} equal to all
of $\Sh$ would settle it.

There is a catch, and it is the reason the problem is hard. Write $\Sh^{\circ}$
for the set of points where the Hodge group of $A_{t}$ is all of
$\SU(V,H)$. Its complement is a countable union of proper special
subvarieties, so $\Sh^{\circ}$ is the complement of a meagre set. Every
point in the Hecke orbit of $s_{0}$ is a CM point, with a torus for its
Hodge group, so it lies \emph{outside} $\Sh^{\circ}$. To close the locus
we need at least one point of $\Sigma\cap\Sh^{\circ}$, and no such point
is known for $n\ge3$ in general.

\section{The first route: strata with big monodromy}

Suppose that we had found one point $t\in\Sigma\cap\Sh^{\circ}$, and let
$S$ be an irreducible stratum of Lemma~\ref{lem:countable} through it. The
first thing to ask is what $S$ can look like. The following lemma says that
$S$ cannot be a small, degenerate piece of the family.

\begin{theorem}\label{thm:monodromy}
Let $S\subseteq\Sh$ be an irreducible closed subvariety of positive
dimension that contains a point of $\Sh^{\circ}$. Then the connected
algebraic monodromy group of the local system $R^{1}\pi_{*}\QQ$ restricted
to the smooth locus of $S$ is the whole group $\SU(V,H)$.
\end{theorem}

\begin{proof}
Let $S^{\mathrm{reg}}$ be the smooth locus of $S$. It is connected, since
$S$ is irreducible. The Hodge group of a very general point of $S$ contains
the Hodge group of every point of $S$, so it contains $\SU(V,H)$. On the
other hand every Hodge group in the family lies in $\SU(V,H)$ \cite{BhH8}.
So the generic Hodge group on $S$ is $\SU(V,H)$.

By Andr\'e's theorem \cite{And92}, the connected algebraic monodromy group
$M$ of a polarisable variation of Hodge structure over a smooth connected
base is a normal subgroup of the derived group of the generic
Mumford--Tate group. Here that derived group is $\SU(V,H)$, which is an
outer form of $\SL_{2n}$ over $\QQ$, and so it is almost simple. Its
connected normal subgroups are $1$ and itself.

If $M=1$, the monodromy is finite. After a finite \'etale cover of
$S^{\mathrm{reg}}$ the local system becomes trivial, and the period map
lifts to a holomorphic map from that cover to $D$. Now $D$ is a bounded
domain, so the coordinate functions of the lift are bounded holomorphic
functions. Pass to a resolution and a smooth compactification of the
cover. By the Riemann extension theorem the bounded functions extend to the
compactification, which is compact and connected, so they are constant.
Then the period map is constant on $S^{\mathrm{reg}}$. But the period map of
the universal family over $\Sh$ is the identity of $\Gamma\backslash D$, so
$S$ is a point, which contradicts $\dim S>0$. Hence $M=\SU(V,H)$.
\end{proof}

So a positive-dimensional stratum through a general point carries all the
monodromy the family has. It is natural to hope that a cycle family with
this much monodromy can be moved off $S$. That hope is exactly the
variational Hodge conjecture for this family, and here is why the
monodromy does not help on its own. The class $\omega$ is invariant under
all of $\SU(V,H)$, so the monodromy never moves $\omega$, and it gives no
constraint that would force the stratum to grow. Big monodromy rules out
degenerate strata. It does not produce new cycles.

\section{The second route: deforming one cycle}

The direct way to move a cycle is deformation theory. Let $Z\subset A_{t}$
be a local complete intersection of codimension $n$. Bloch \cite{Blo72}
defined the \emph{semiregularity map}
\[
  \sigma_{Z}\colon H^{1}(Z,N_{Z})\longrightarrow H^{n+1}(A_{t},\Omega^{n-1}),
\]
and proved the following. If $\sigma_{Z}$ is injective, then $Z$ deforms
along every deformation of $A_{t}$ along which its class stays of Hodge type.
Buchweitz and Flenner \cite{BF03} extended the definition and the theorem
to perfect complexes, and Pridham \cite{Pri24} gave a general proof.

\begin{theorem}\label{thm:semireg}
Suppose that at some point $t\in\Sh$ there is a local complete intersection
$Z\subset A_{t}$ of codimension $n$ whose class is
$\cl(Z)=\alpha+c\,\eta^{n}$ with $\alpha$ a nonzero Weil class and
$c\in\QQ$, and whose semiregularity map is injective. Then $\Sigma=\Sh$.
\end{theorem}

\begin{proof}
Along the universal family over a small neighbourhood $U$ of $t$, the flat
transport of $\cl(Z)$ is $\alpha+c\eta^{n}$, and both terms are Hodge
classes at every point, because the Weil plane and $\eta$ are. So the class
of $Z$ stays of Hodge type along every direction of $U$. By Bloch's theorem,
$Z$ deforms to every order along every curve germ in $U$. Then the image of
its Hilbert scheme component contains the formal neighbourhood of $t$. Since
that image is a closed analytic subset of $U$, it contains a neighbourhood
of $t$. So one stratum of Lemma~\ref{lem:countable} has nonempty interior,
and it is all of $\Sh$ by irreducibility. At every point $s$ we then have an
algebraic cycle of class $\alpha(s)+c\eta^{n}$, so $\alpha(s)$ is
algebraic.

It remains to pass from $\alpha$ to $\omega$. The field $K$ acts on the
Weil plane: an element $\tau\in K^{\times}$ acts on $\bigwedge^{2n}_{K}V$ by
multiplication by $\tau^{2n}$. A positive integer multiple $m\tau$ is a
genuine endomorphism of $A_{s}$. Its pullback, divided by $m^{2n}$, is
$\tau^{*}$ on $H^{2n}$, so $\tau^{*}$ carries algebraic classes to
algebraic classes. Choose $\tau$ whose argument is not a rational multiple of $\pi$, for
example $\tau=a+\sqrt{-d}$ with $a$ a positive integer chosen so that
$\tau/\bar\tau$ is not a root of unity. Then $\tau^{2n}\notin\QQ$, and the classes
$\alpha$ and $\tau^{*}\alpha$ span the Weil plane over $\QQ$, so $\omega(s)$
is algebraic.
\end{proof}

The same holds with $Z$ replaced by a perfect complex $E$ on $A_{t}$, with
the Chern character in place of the cycle class, by the theorem of
Buchweitz--Flenner and Pridham. This is the route Markman followed for
sixfolds of split Weil type, using secant sheaves and an equivariant form
of semiregularity \cite{Mar25b}.

\section{Where the attempt stops}\label{sec:stops}

Both routes need an input that we do not have.

\subsection*{The first route needs a point of $\Sigma$ with full Hodge group.}
Outside the families where the conjecture is already known, such as
Markman's split sixfolds, every point at which $\omega$ is known to be
algebraic is a CM point or lies on a proper special subvariety. These points are excluded from $\Sh^{\circ}$ by definition.
Theorem~\ref{thm:monodromy} tells us what a stratum through a point of
$\Sh^{\circ}$ would look like, but it gives no way to find one.

\subsection*{The second route needs a semiregular object, and the obvious
ones are not.}
For an object whose Chern character is \emph{exactly} a Weil class, with no
other component, the semiregularity map is never injective, at any member
of any Weil family (\cite{BhH8}, the theorem labelled \texttt{thm:p2false}
in its source). The reason is
Voisin's. Such an object would deform along every $K$-linear deformation of
$A_{t}$, including the non-algebraic complex tori of Weil type. A very
general such torus has no analytic subvarieties of intermediate dimension,
and its coherent sheaves have no Chern classes in intermediate degree
\cite{Voi02b}. So the deformation cannot exist, and the semiregularity map
cannot be injective. With a polarisation term allowed, as in
Theorem~\ref{thm:semireg}, the criterion survives, but it asks for an object
with a precise and large amount of self-extension. For $n\ge3$ and a
general CM field no such object is known. A list of the natural
constructions that provably fail is in \cite{BhH8}.

\subsection*{Neither gap is a technicality.}
If $\Sigma=\Sh$, then $\Sigma$ meets $\Sh^{\circ}$, so the first gap is
a consequence of the goal as well as a step towards it. And by
\cite{BhH8} (the corollary labelled \texttt{cor:weilequiv}), closing the
locus is equivalent to one case of the Lefschetz standard conjecture: the
case of the total space of the family over a projective curve through
$s_{0}$ and a very general point. These are the variational Hodge
conjecture and the standard conjectures in disguise. They are not small residuals.

\section{What this note establishes}

Proved here, or proved in the cited work and used as theorems:
\begin{enumerate}[label=(\arabic*)]
\item The algebraic locus is a countable union of closed algebraic
subsets, it is dense, and it is either everything or meagre.
\item Any positive-dimensional stratum of it through a member with full
Hodge group has full monodromy $\SU(V,H)$.
\item A single semiregular cycle or complex, at any member, with class a
nonzero Weil class plus a multiple of $\eta^{n}$, would close the locus,
and with it the Hodge conjecture for every abelian variety of Weil type in
that family.
\end{enumerate}
Not proved: that such a cycle or complex exists, or that $\Sigma$ contains
a single member with full Hodge group, when $n\ge3$. So the Hodge
conjecture for abelian varieties of Weil type of dimension at least six
remains open, and with it the Hodge conjecture.

\begin{thebibliography}{99}

\bibitem{And92} Y.~Andr\'e, \emph{Mumford--Tate groups of mixed Hodge
structures and the theorem of the fixed part}, Compositio Math.\
\textbf{82} (1992), 1--24.

\bibitem{BB66} W.~L.~Baily, Jr.\ and A.~Borel, \emph{Compactification of
arithmetic quotients of bounded symmetric domains}, Ann.\ of Math.\ (2)
\textbf{84} (1966), 442--528.

\bibitem{BhH8} D.~Bhattacharjee, \emph{Algebraic Loci of Weil Classes on
Abelian Varieties}, LaTeX source, verification code and Lean certificate,
release v5.1.0, \url{https://github.com/creelie/H8}; archived at
Zenodo, \url{https://doi.org/10.5281/zenodo.22950275}.

\bibitem{Blo72} S.~Bloch, \emph{Semi-regularity and de Rham cohomology},
Invent.\ Math.\ \textbf{17} (1972), 51--66.

\bibitem{BF03} R.-O.~Buchweitz and H.~Flenner, \emph{A semiregularity map
for modules and applications to deformations}, Compositio Math.\
\textbf{137} (2003), 135--210.

\bibitem{Mar25a} E.~Markman, \emph{Cycles on abelian $2n$-folds of Weil
type from secant sheaves on abelian $n$-folds}, preprint,
arXiv:2502.03415 (2025).

\bibitem{Mar25b} E.~Markman, \emph{Secant sheaves and Weil classes on
abelian varieties}, preprint, arXiv:2509.23403 (2025).

\bibitem{MZ99} B.~Moonen and Yu.~G.~Zarhin, \emph{Hodge classes on abelian
varieties of low dimension}, Math.\ Ann.\ \textbf{315} (1999), 711--733.

\bibitem{Pri24} J.~P.~Pridham, \emph{Semiregularity as a consequence of
Goodwillie's theorem}, Forum Math.\ Sigma \textbf{12} (2024), e126.

\bibitem{Voi02b} C.~Voisin, \emph{A counterexample to the Hodge conjecture
extended to K\"ahler varieties}, Int.\ Math.\ Res.\ Not.\ \textbf{2002},
no.~20, 1057--1075.

\bibitem{Weil77} A.~Weil, \emph{Abelian varieties and the Hodge ring},
in: \OE uvres Scientifiques, Vol.~III, Springer, New York, 1979,
421--429.

\end{thebibliography}

\end{document}
